
\documentclass[12pt]{article}
\usepackage{amsmath, amssymb}
\usepackage{geometry}
\geometry{margin=1in}
\title{Scalar--Angular Theory (SAT): Lay Overview}
\author{}
\date{}

\begin{document}

\maketitle

\section*{For Readers Curious About the Big Picture (Not the Equations)}

For over a century, two towering theories have formed the backbone of our understanding of the universe:
\begin{itemize}
  \item \textbf{General Relativity}, which describes how gravity works and how space and time bend;
  \item \textbf{Quantum Field Theory}, which describes the behavior of particles and forces at the smallest scales.
\end{itemize}

Each theory performs remarkably well in its own domain. But when we try to bring them together---to describe, for instance, black holes or the very early universe---they don’t quite align. They behave like puzzle pieces cut from different sets.

Our project, which we call \textit{Scalar--Angular Theory} (SAT), doesn’t seek to replace those puzzle pieces. Instead, we aim to build a frame that might help both fit together more naturally. One could think of SAT as scaffolding: not the building itself, but a temporary structure that allows us to make sense of complex parts while still under construction.

We start with very simple elements:
\begin{itemize}
  \item \textbf{Filaments}: lines that stretch through space and time;
  \item \textbf{Wavefronts}: sweeping three-dimensional surfaces that move and interact with those filaments.
\end{itemize}

From these minimal ingredients, SAT attempts to show how features like particles, gravity, and the fabric of space itself could emerge---simply by studying how filaments are arranged, how they twist, and how they intersect with the wavefront.

This work is ongoing. Some of the ideas remain rough; portions of the mathematics were assisted by tools such as ChatGPT. But our hope is that, even if SAT turns out not to be the final answer, it might help us reinterpret known theories in a new light---and perhaps bring them a little closer together.

\section*{What SAT Is Not}

SAT is not a replacement for General Relativity or Quantum Field Theory. It is not a claim to final truth or a complete theory of everything. It does not attempt to replicate existing models with different jargon, nor does it insist that all known physics be reducible to its terms.

We do not assert SAT as “more correct” than current theories. Instead, we explore it as a scaffold---a geometric and conceptual framework that might assist in revealing how known physics could fit together more coherently.

\end{document}
